%=====================================================================
\section{Token-Level Analysis: What an Embedding Edit Removes}
\label{sec:token}
%=====================================================================

A natural concern with any embedding edit is whether it removes the concept from a word or the word from the model.
Since the input embedding is where the model first reads a token, we expect an edit to control \emph{access} to the concept through the edited words rather than the knowledge itself, and we expect stronger erasure to cost word integrity.
To test this, we probe every edited word with the integrity tests of \S\ref{sec:setup} and compare each method with two references: the word deleted, and the word replaced by a nonce.

\paragraph{Embedding erasure blocks access}
On efficacy questions that contain none of the edited tokens, every edited model answers exactly like the unedited one (efficacy \expnum{0.000} in all \expnum{18} configurations checked, including $r = 100$).
The concept's knowledge thus remains in the weights and is only cut off from the edited words.
At the operating points selected by our rule, efficacy lies close to the floor that deleting the edited words from the questions would produce.

\paragraph{Removing more costs the words}
At these operating points, about half of the edited words can no longer be read (\expnum{0.49--0.61}), and most of the concept-exclusive ones (\expnum{0.65--0.81}), with no significant difference between methods.
Requiring the words to remain copyable (a post hoc rule) caps efficacy at \expnum{0.06--0.14}; under this constraint \cinder{} removes about twice as much as \ember{}, whereas the advantage of \cinderz{} disappears, since part of it came from damage to the surrounding prompt, consistent with its extra component being a direction that every token shares (\S\ref{sec:cinderz}).
Table~\ref{tab:examples} shows what destruction looks like in free text.

\paragraph{A nonce-like edit is only partly reachable}
An erased concept word should behave like a nonce: copyable, but meaningless.
A learned edit (post hoc) that pushes each edited word toward the behaviour of its nonce keeps the words copyable and removes more than any subtraction at readable strength (\expnum{2--3$\times$} \ember{}'s, \expnum{1.3--1.6$\times$} \cinder{}'s), but it reaches a substantial nonce floor in only \expnum{2} of 12 concepts, and the model often can no longer name the word (\S\ref{app:nonce}).
This suggests that, within a single embedding row, the concept and the word's identity share the same few directions.
\expfill{Figure: efficacy against word integrity (copy) along the $r$ grid for each method, with the nonce and deletion floors.}

\begin{table}[t]
  \centering
  \small
  \setlength{\tabcolsep}{3pt}
  \begin{tabular}{@{}P{2.0cm}P{2.45cm}P{2.45cm}@{}}
    \toprule
    Prompt & $\delta = 1$ ($r \approx 0.6$) & $\delta = 200$ (selected by $H$) \\
    \midrule
    Who is Prince \textbf{Harry}, and who is his wife? & \textit{Prince Harry is the second son of King Charles III [\dots] His wife is Meghan Markle} & \textit{Prince was an American singer, songwriter [\dots] Prince's wife was Mayte} \\
    \addlinespace
    What country does the toy company \textbf{Lego} come from [\dots]? & \textit{Lego comes from Denmark.} & \textit{Please provide me with the name of the toy company!} \\
    \addlinespace
    In which year was Donald \textbf{Trump} first elected President? & \textit{[\dots] in 2016.} & \textit{Donald Trump} \\
    \bottomrule
  \end{tabular}
  \caption{Responses of Gemma-2-2B-it after erasing Harry Potter with \ember{} (seed 42) at a moderate strength and at the strength that \ember{}'s selection score chooses. Edited tokens are in bold; all three belong to the Harry Potter feature. \note{verify $r$ at $\delta = 1$ and that the selection chose $\delta = 200$ for this run.}}
  \label{tab:examples}
\end{table}
